% GENERATED FILE. Edit canonical lessons, then run npm run generate:books.
% canonical-source-hash: fnv1a64:6944da716e4fff24
% canonical-lessons: LA-C307-venenum, LA-C307-veneficus, LA-C307-insidiae, LA-C307-coniuratio, LA-C307-proditor

\chapter{Venēnum}
\label{ch:la-venenum}
\hlchaptermodality{307}

\begin{chapteropening}
\textbf{By the end of this chapter:} \emph{I can talk about poison and plots: poison, a poisoner, an ambush, a conspiracy and a traitor.}

Complete the last lesson of chapter 307: I can talk about poison and plots: poison, a poisoner, an ambush, a conspiracy and a traitor.
\end{chapteropening}

\section[venēnum, venēnī]{venēnum, venēnī — poison; a drug}

\label{lesson:LA-C307-venenum}



\begin{quote}
Before the new one: say the Latin for a dagger, then the Latin for to kill.
\end{quote}

\subsection*{You'll want to know: venēnum, venēnī}
\textbf{venēnum, venēnī} — \textquotedblleft{}poison; a drug\textquotedblright{}.

\begin{quote}
Medicus venēnum in pōculō invēnit.
\end{quote}

\begin{quote}
\emph{The doctor found poison in the cup.}
\end{quote}

Neuter. First any strong drug or charm, then poison. English \emph{venom} comes from it through French.

\begin{cousinweb}
\textbf{venēnum} is inherited; EtymDB traces it to Proto-Italic \emph{*wenesnom}, from a root meaning \textquotedblleft{}to desire\textquotedblright{}, so first a love potion. English \emph{venom} comes from it through French.
\end{cousinweb}

\begin{grammarlens}[title={What you've built}]
One more word for telling, explaining and describing at B1.
\end{grammarlens}

\subsection*{Your turn}
\begin{itemize}
\raggedright
  \item \emph{Say it:} \emph{venēnum, venēnī}
  \item \emph{Say it:} \emph{venēnum, venēnī}, once more
  \item \emph{Say it:} a sentence of your own with \emph{venēnum, venēnī}
  \item \emph{Recall:} say the Latin for a dagger, then the Latin for to kill, then say \emph{venēnum, venēnī} again
\end{itemize}

\begin{tcolorbox}[breakable,colback=teal!4,colframe=teal!35!black,title={Before you move on}]
What does venēnum, venēnī mean? (\textquotedblleft{}Poison; a drug\textquotedblright{}.) Say it once more.
\end{tcolorbox}
\section[venēficus, venēficī]{venēficus, venēficī — a poisoner; a sorcerer}

\label{lesson:LA-C307-veneficus}



\begin{quote}
Before the new one: say the Latin for to kill, then the Latin for poison.
\end{quote}

\subsection*{You'll want to know: venēficus, venēficī}
\textbf{venēficus, venēficī} — \textquotedblleft{}a poisoner; a sorcerer\textquotedblright{}.

\begin{quote}
Venēficus herbās perīculōsās in hortō alēbat.
\end{quote}

\begin{quote}
\emph{The poisoner grew dangerous herbs in his garden.}
\end{quote}

Masculine; the feminine \emph{venēfica} is the witch of Roman stories.

\begin{cousinweb}
\textbf{venēficus} joins \emph{venēnum}, \textquotedblleft{}poison\textquotedblright{}, and \emph{faciō}, \textquotedblleft{}to make\textquotedblright{} (EtymDB, Lewis \& Short): a maker of poisons and charms.
\end{cousinweb}

\begin{grammarlens}[title={What you've built}]
One more word for telling, explaining and describing at B1.
\end{grammarlens}

\subsection*{Your turn}
\begin{itemize}
\raggedright
  \item \emph{Say it:} \emph{venēficus, venēficī}
  \item \emph{Say it:} \emph{venēficus, venēficī}, once more
  \item \emph{Say it:} a sentence of your own with \emph{venēficus, venēficī}
  \item \emph{Recall:} say the Latin for to kill, then the Latin for poison, then say \emph{venēficus, venēficī} again
\end{itemize}

\begin{tcolorbox}[breakable,colback=teal!4,colframe=teal!35!black,title={Before you move on}]
What does venēficus, venēficī mean? (\textquotedblleft{}A poisoner; a sorcerer\textquotedblright{}.) Say it once more.
\end{tcolorbox}
\section[īnsidiae, īnsidiārum]{īnsidiae, īnsidiārum — an ambush; a plot}

\label{lesson:LA-C307-insidiae}



\begin{quote}
Before the new one: say the Latin for poison, then the Latin for a poisoner.
\end{quote}

\subsection*{You'll want to know: īnsidiae, īnsidiārum}
\textbf{īnsidiae, īnsidiārum} — \textquotedblleft{}an ambush; a plot\textquotedblright{}.

\begin{quote}
Hostēs in silvā īnsidiās parāvērunt.
\end{quote}

\begin{quote}
\emph{The enemy prepared an ambush in the forest.}
\end{quote}

Feminine, plural only: men sitting hidden, waiting. English \emph{insidious}.

\begin{cousinweb}
\textbf{īnsidiae} is built on \emph{īnsideō}, \textquotedblleft{}to sit in, to occupy\textquotedblright{} (Lewis \& Short), from \emph{in} and \emph{sedeō}, \textquotedblleft{}to sit\textquotedblright{}: a lying in wait. English \emph{insidious} comes from it.
\end{cousinweb}

\begin{grammarlens}[title={What you've built}]
One more word for telling, explaining and describing at B1.
\end{grammarlens}

\subsection*{Your turn}
\begin{itemize}
\raggedright
  \item \emph{Say it:} \emph{īnsidiae, īnsidiārum}
  \item \emph{Say it:} \emph{īnsidiae, īnsidiārum}, once more
  \item \emph{Say it:} a sentence of your own with \emph{īnsidiae, īnsidiārum}
  \item \emph{Recall:} say the Latin for poison, then the Latin for a poisoner, then say \emph{īnsidiae, īnsidiārum} again
\end{itemize}

\begin{tcolorbox}[breakable,colback=teal!4,colframe=teal!35!black,title={Before you move on}]
What does īnsidiae, īnsidiārum mean? (\textquotedblleft{}An ambush; a plot\textquotedblright{}.) Say it once more.
\end{tcolorbox}
\section[coniūrātiō, coniūrātiōnis]{coniūrātiō, coniūrātiōnis — a conspiracy}

\label{lesson:LA-C307-coniuratio}



\begin{quote}
Before the new one: say the Latin for a poisoner, then the Latin for an ambush.
\end{quote}

\subsection*{You'll want to know: coniūrātiō, coniūrātiōnis}
\textbf{coniūrātiō, coniūrātiōnis} — \textquotedblleft{}a conspiracy\textquotedblright{}.

\begin{quote}
Cōnsul coniūrātiōnem in senātū patefēcit.
\end{quote}

\begin{quote}
\emph{The consul exposed the conspiracy in the senate.}
\end{quote}

Feminine, 3rd declension. Cicero brought one to light in 63 BC.

\begin{cousinweb}
\textbf{coniūrātiō} is built on \emph{coniūrō}, \textquotedblleft{}to swear together\textquotedblright{}, from \emph{con-} and \emph{iūrō}, \textquotedblleft{}to swear\textquotedblright{} (Lewis \& Short): a band bound by a shared oath.
\end{cousinweb}

\begin{grammarlens}[title={What you've built}]
One more word for telling, explaining and describing at B1.
\end{grammarlens}

\subsection*{Your turn}
\begin{itemize}
\raggedright
  \item \emph{Say it:} \emph{coniūrātiō, coniūrātiōnis}
  \item \emph{Say it:} \emph{coniūrātiō, coniūrātiōnis}, once more
  \item \emph{Say it:} a sentence of your own with \emph{coniūrātiō, coniūrātiōnis}
  \item \emph{Recall:} say the Latin for a poisoner, then the Latin for an ambush, then say \emph{coniūrātiō, coniūrātiōnis} again
\end{itemize}

\begin{tcolorbox}[breakable,colback=teal!4,colframe=teal!35!black,title={Before you move on}]
What does coniūrātiō, coniūrātiōnis mean? (\textquotedblleft{}A conspiracy\textquotedblright{}.) Say it once more.
\end{tcolorbox}
\section[prōditor, prōditōris]{prōditor, prōditōris — a traitor}

\label{lesson:LA-C307-proditor}



\begin{quote}
Before the new one: say the Latin for an ambush, then the Latin for a conspiracy.
\end{quote}

\subsection*{You'll want to know: prōditor, prōditōris}
\textbf{prōditor, prōditōris} — \textquotedblleft{}a traitor\textquotedblright{}.

\begin{quote}
Prōditor portās hostibus aperuit.
\end{quote}

\begin{quote}
\emph{The traitor opened the gates to the enemy.}
\end{quote}

Masculine, 3rd declension, from \emph{prōdō}, 'to betray'. \emph{Prōditor patriae}: a traitor to his country.

\begin{cousinweb}
\textbf{prōditor} is built on \emph{prōdō}, \textquotedblleft{}to give up, to betray\textquotedblright{} (Lewis \& Short), with the ending \emph{-tor}.
\end{cousinweb}

\begin{grammarlens}[title={What you've built}]
One more word for telling, explaining and describing at B1.
\end{grammarlens}

\subsection*{Your turn}
\begin{itemize}
\raggedright
  \item \emph{Say it:} \emph{prōditor, prōditōris}
  \item \emph{Say it:} \emph{prōditor, prōditōris}, once more
  \item \emph{Say it:} a sentence of your own with \emph{prōditor, prōditōris}
  \item \emph{Recall:} say the Latin for an ambush, then the Latin for a conspiracy, then say \emph{prōditor, prōditōris} again
\end{itemize}

\begin{tcolorbox}[breakable,colback=teal!4,colframe=teal!35!black,title={Before you move on}]
What does prōditor, prōditōris mean? (\textquotedblleft{}A traitor\textquotedblright{}.) Say it once more.
\end{tcolorbox}
